\section{Weitere reflexive Verben (nach Lingolia)}
\label{sec:WeitereReflexiv}

Vollständige Liste der wichtigsten reflexiven Verben nach Niveau:

\subsection*{A1/A2 - reflexive Verben mit Akkusativ}
\begin{center}
\begin{tabular}{|l|c|l|}
\hline
Verb & Reflexivpronomen & Beispiel \\ \hline
sich duschen & mich & Ich dusche mich. \\ \hline
sich waschen & mich & Du wäschst dich. \\ \hline
sich anziehen & mich & Ich ziehe mich an. \\ \hline
sich ausziehen & mich & Er zieht sich aus. \\ \hline
sich baden & mich & Sonntags bade ich mich. \\ \hline
sich befinden & sich & Das Hotel befindet sich in Berlin. \\ \hline
sich beschweren & sich & Ich beschwere mich. \\ \hline
sich freuen & mich & Wir freuen uns auf den Urlaub. \\ \hline
sich fühlen & mich & Ich fühle mich gut. \\ \hline
sich informieren & mich & Sie informiert sich. \\ \hline
sich kämmen & mich & Kämm dich! \\ \hline
sich legen & mich & Ich lege mich hin. \\ \hline
sich rasieren & mich & Er rasiert sich. \\ \hline
sich setzen & mich & Setzen Sie sich! \\ \hline
sich stellen & mich & Er stellt sich an die Tür. \\ \hline
sich umziehen & mich & Ich ziehe mich um. \\ \hline
sich verlaufen & mich & Ich habe mich verlaufen. \\ \hline
sich erholen & mich & Wir erholen uns gut. \\ \hline
sich beeilen & mich & Beeil dich! \\ \hline
sich entscheiden & sich & Ich entscheide mich. \\ \hline
\end{tabular}
\end{center}

\subsection*{A1/A2 - reflexive Verben mit Dativ}
\begin{center}
\begin{tabular}{|l|c|l|}
\hline
Verb & Reflexivpronomen & Beispiel \\ \hline
sich die Haare bürsten & mir & Du bürstest dir die Haare. \\ \hline
sich die Haare kämmen & mir & Ich kämme mir die Haare. \\ \hline
sich etwas anziehen & mir & Ich ziehe mir eine Jacke an. \\ \hline
sich etwas leisten & mir & Das kann ich mir nicht leisten. \\ \hline
sich etwas vorstellen & mir & Ich stelle mir das vor. \\ \hline
sich die Zähne putzen & mir & Ich putze mir die Zähne. \\ \hline
sich die Hände waschen & mir & Ich wasche mir die Hände. \\ \hline
\end{tabular}
\end{center}

\section{Reflexive Verben mit festen Präpositionen}
\label{sec:ReflexivMitPraep}

\begin{center}
\begin{tabular}{|l|c|c|c|}
\hline
Verb & Präposition & Beispiel & Bedeutung \\ \hline
sich ärgern & über & Ich ärgere mich über dich. & be annoyed at \\ \hline
sich freuen & auf & Ich freue mich auf das Fest. & look forward to \\ \hline
sich freuen & über & Ich freue mich über das Geschenk. & be happy about \\ \hline
sich interessieren & für & Ich interessiere mich für Musik. & be interested in \\ \hline
sich entscheiden & für & Er entscheidet sich für das Angebot. & decide for \\ \hline
sich beschweren & über & Die Kunden beschweren sich über den Service. & complain about \\ \hline
sich vorbereiten & auf & Wir bereiten uns auf die Prüfung vor. & prepare for \\ \hline
sich gewöhnen & an & Ich gewöhne mich an das Klima. & get used to \\ \hline
sich erinnern & an & Ich erinnere mich an dich. & remember \\ \hline
sich beteiligen & an & Wir beteiligen uns an dem Projekt. & participate in \\ \hline
sich wenden & an & Ich wende mich an den Chef. & turn to \\ \hline
sich verabschieden & von & Ich verabschiede mich von dir. & say goodbye \\ \hline
sich verlieben & in & Er hat sich in sie verliebt. & fall in love \\ \hline
sich unterscheiden & von & Das unterscheidet sich von dem anderen. & differ from \\ \hline
sich verstehen & mit & Ich verstehe mich gut mit ihm. & get along with \\ \hline
sich bewerben & um & Ich bewerbe mich um die Stelle. & apply for \\ \hline
sich kümmern & um & Sie kümmert sich um das Kind. & take care of \\ \hline
\end{tabular}
\end{center}

\section{Übungen zu reflexiven Verben}
\label{sec:ReflexivUebungen}

\subsection*{Übung 1: Setze das richtige Reflexivpronomen ein}
\begin{enumerate}
    \item Ich wasche \underline{\hspace{1cm}} (mich/mir) die Hände.
    \item Er rasiert \underline{\hspace{1cm}} (sich/sich) jeden Morgen.
    \item Wir freuen \underline{\hspace{1cm}} (uns/euch) auf die Reise.
    \item Sie kämmt \underline{\hspace{1cm}} (sich/sich) die Haare.
    \item Ich ziehe \underline{\hspace{1cm}} (mich/mir) eine Jacke an.
    \item Die Kinder erinnern \underline{\hspace{1cm}} (sich/sich) an den Film.
    \item Wir bereiten \underline{\hspace{1cm}} (uns/euch) auf die Prüfung vor.
    \item Er interessiert \underline{\hspace{1cm}} (sich/sich) für Sport.
\end{enumerate}

\loesung{
\begin{enumerate}
    \item mir (Dativ - weil "die Hände" das Akkusativobjekt ist)
    \item sich
    \item uns
    \item sich (Dativ - "die Haare" ist Akkusativobjekt)
    \item mir (Dativ - "eine Jacke" ist Akkusativobjekt)
    \item sich (Akkusativ)
    \item uns (Akkusativ)
    \item sich (Akkusativ)
\end{enumerate}
}

\chapterend